\section{Related Work}
\textit{Image-only description and evaluation.} Reference-free metrics such as CLIPScore score the compatibility between an image and a caption \cite{hessel2021}, and captioning datasets for BLV users, such as VizWiz-Captions, target general scene description \cite{gurari2020}. None of these reason about what the surrounding page already tells the reader. Kreiss et al.\ showed that referenceless metrics align poorly with accessibility judgments precisely because context is missing, and called for context-aware evaluation \cite{kreiss2022}.

\textit{Context-aware generation.} Stangl et al.\ found that BLV users' information wants depend on both the image source and the user's goal \cite{stangl2021}. Mohanbabu and Pavel (ASSETS 2024) extract webpage context and combine it with the image to produce context-aware descriptions; in their BLV user study, these were generally preferred to context-free descriptions \cite{mohanbabu2024}. We use this system, hereafter \base, as our baseline. It is, however, general-purpose rather than tuned to product evaluation. In our reproduction on Amazon.fr listings, it often produced details that were visually correct but irrelevant to the purchase, such as a model's jeans, necklaces, or handbag. Such details inflate apparent novelty without helping the shopper. \sys therefore treats relevance to the target product and redundancy with the DOM as explicit objectives.
